% ***************************************************
% Appendix
% ***************************************************

\chapter{Example \texttt{AOFLAGGER} Lua script}
\label{app:aoflagger}
The \texttt{AOFLAGGER} strategy used for the flagging of RFI in this thesis consists of two targeted passes, and a final general pass, designed to address the different RFI environments present across the full 8--12 GHz frequency band. An example of the script is provided in this appendix, the complete scripts used during the calibration process are available in the associated GitHub repository under `Data Processing' and the relevant observation ID.\footnote{\url{https://github.com/MoonlitMoo/Masters-2025}}

The \texttt{options} function defines the spectral windows that are flagged by each pass, and the corresponding execution function. The first pass (\texttt{execute\_awr}) targets spectral windows 10--11, corresponding to frequencies $\approx$9.3--9.5 GHz which contain intermittent satellite RFI. The second pass (\texttt{execute\_microwave}) covers spectral windows 22--31 (approximately 10.7--12 GHz), where continuous ground-based communication interference is present. The final pass (\texttt{execute\_final}) applies to the full frequency band 8--12 GHz. The targeted passes use independent lower flagging thresholds than the final pass, providing adaptability to the amount of RFI measured in an observation. The values shown are representative of those used in the calibration process, with thresholds ranging from 0.65 to 5.0.

\begin{lstlisting}[
    language=Lua, 
    caption={Example \texttt{AOFLAGGER} Lua script.}, 
    label={lst:reduction},
    literate={é}{{\'e}}1, 
    numberstyle=\tiny,
    basicstyle=\ttfamily\footnotesize
]
aoflagger.require_min_version("3.0")

function options()
  opts1 = {
    ["bands"] = {10,11},
    ["execute-function"] = "execute_awr",

  }

  opts2 = {
    ["bands"] = {22,23,24,25,26,27,28,29,30,31},
    ["execute-function"] = "execute_microwave"
  }

  opts3 = {
    ["execute-function"] = "execute_final"
  }

  return {
    ["1 - weather radars"] = opts1,
    ["2 - microwave link"] = opts2,
    ["3 - final sweep"] = opts3
  }
end

function execute_awr(input)
    -- 9.3-9.5 GHz (airborne weather radars)
    trim_rfi(input, 1.0, 1.0)
end

function execute_microwave(input)
    -- 10.7-12 GHz (Microwave link + geostationary satellite near 11.7-12 GHz)
    trim_rfi(input, 1.0, 1.0)
end

function execute_final(input)
    -- Run over remaining at the end as cleanup
    trim_rfi(input, 2.0, 2.0)
end

function trim_rfi(input, base_threshold, transient_threshold_factor)
  -- base_threshold, lower means more sensitive detection
  -- transient_threshold_factor, decreasing this value makes detection of transient RFI more aggressive
  local flag_polarizations = input:get_polarizations()
  local flag_representations = { "amplitude" }
  local iteration_count = 4 -- how many iterations to perform?
  local threshold_factor_step = 1.75 -- How much to increase the sensitivity each iteration?
  local exclude_original_flags = true
  local frequency_resize_factor = 1.0 -- Amount of "extra" smoothing in frequency direction

  local inpPolarizations = input:get_polarizations()

  if not exclude_original_flags then
    input:clear_mask()
  end
  local copy_of_input = input:copy()

  for ipol, polarization in ipairs(flag_polarizations) do
    local pol_data = input:convert_to_polarization(polarization)
    local converted_data
    local converted_copy

    for _, representation in ipairs(flag_representations) do
      converted_data = pol_data:convert_to_complex(representation)
      converted_copy = converted_data:copy()

      for i = 1, iteration_count - 1 do
        local threshold_factor = threshold_factor_step ^ (iteration_count - i)

        local sumthr_level = threshold_factor * base_threshold
        if exclude_original_flags then
          aoflagger.sumthreshold_masked(
            converted_data,
            converted_copy,
            sumthr_level,
            sumthr_level * transient_threshold_factor,
            true,
            true
          )
        else
          aoflagger.sumthreshold(converted_data, sumthr_level, sumthr_level * transient_threshold_factor, true, true)
        end

        -- Do timestep & channel flagging
        local chdata = converted_data:copy()
        aoflagger.threshold_timestep_rms(converted_data, 2.0)
        aoflagger.threshold_channel_rms(chdata, 1.5 * threshold_factor, true)
        converted_data:join_mask(chdata)

        -- High pass filtering steps
        converted_data:set_visibilities(converted_copy)
        if exclude_original_flags then
          converted_data:join_mask(converted_copy)
        end

        local resized_data = aoflagger.downsample(converted_data, 3, frequency_resize_factor, true)
        aoflagger.low_pass_filter(resized_data, 21, 31, 2.5, 5.0)
        aoflagger.upsample(resized_data, converted_data, 3, frequency_resize_factor)

        local tmp = converted_copy - converted_data
        tmp:set_mask(converted_data)
        converted_data = tmp

        aoflagger.visualize(converted_data, "Residual #" .. i, i + iteration_count)
        aoflagger.set_progress((ipol - 1) * iteration_count + i, #flag_polarizations * iteration_count)
      end -- end of iterations

      if exclude_original_flags then
        aoflagger.sumthreshold_masked(
          converted_data,
          converted_copy,
          base_threshold,
          base_threshold * transient_threshold_factor,
          true,
          true
        )
      else
        aoflagger.sumthreshold(converted_data, base_threshold, base_threshold * transient_threshold_factor, true, true)
      end
    end -- end of complex representation iteration

    if exclude_original_flags then
      converted_data:join_mask(converted_copy)
    end

    -- Helper function used below
    function contains(arr, val)
      for _, v in ipairs(arr) do
        if v == val then
          return true
        end
      end
      return false
    end

    if contains(inpPolarizations, polarization) then
      if input:is_complex() then
        converted_data = converted_data:convert_to_complex("complex")
      end
      input:set_polarization_data(polarization, converted_data)
    else
      input:join_mask(converted_data)
    end

    aoflagger.visualize(converted_data, "Residual #" .. iteration_count, 2 * iteration_count)
    aoflagger.set_progress(ipol, #flag_polarizations)
  end -- end of polarization iterations

  if exclude_original_flags then
    aoflagger.scale_invariant_rank_operator_masked(input, copy_of_input, 0.2, 0.2)
  else
    aoflagger.scale_invariant_rank_operator(input, 0.2, 0.2)
  end

  aoflagger.threshold_timestep_rms(input, 4.0)

  input:flag_nans()
end
\end{lstlisting}
